\chapter{Практика}

Разработаем pipeline для предобработки данных.

\section{Считывание данных}
Считывание данных проивзодится с помощью библиотеки Pandas. Она позволяет быстро
и удобно считать данные, а также хранит считанные данные в виде таблицы. \\

\section{Первичная визуализация}
Также Pandas позволяет сделать визуализацию для первичного анализа

\begin{figure}[H]
    \centering
    \includegraphics{vis.png}
    \caption{Первичная визуализация}
\end{figure}

\section{Заполнение пропусков}

Модель поддерживает следующие способы заполнения пропусков:

\begin{enumerate}
\item drop\_rows - удалить строку (самый простой вопрос)
\item median - заменить медианной по столбцу
\item meam - заменить средним
\end{enumerate}

Также можно вывести heatmap по пропускам. В моем датасете нет пропусков, поэтому цвет heatmap однотонный:

\begin{figure}[H]
    \centering
    \includegraphics{heatmap.png}
    \caption{Heatpmap по пропускам}
\end{figure}

\section{Распределение целевой переменной}
В программе можно задать целевую переменную (столбец) и вывести ее распределение на графике:

\begin{figure}[H]
    \centering
    \includegraphics{distr.png}
    \caption{Распределение целевой переменной}
\end{figure}

Здесь по оси Х - номера континентов, по оси Y - кол-во стран на этом континенте.



\endinput